% RESEARCH_DRAFT_SHA256: 908aeb1adbda7fa1c89f873c867392e95e1f47aec94a470702fe0ee88bf0d752
% Layout and source-reference wrappers expanded for the web renderer.
\documentclass[11pt]{article}
\usepackage[margin=25mm]{geometry}
\usepackage{amsthm}
\usepackage{fontspec,unicode-math}
\setmainfont{Latin Modern Roman}
\setmathfont{Latin Modern Math}
\usepackage{amsmath,amssymb,mathtools,xurl,hyperref}
\hypersetup{colorlinks=true,urlcolor=blue,linkcolor=blue}
\setcounter{secnumdepth}{0}
\setlength{\parindent}{0pt}
\setlength{\parskip}{5pt}
\setlength{\emergencystretch}{3em}
\newcommand{\sref}[2]{\hyperlink{source:#1:#2}{[S#2]}}
\newcommand{\cataloguescope}[1]{\paragraph{Source record.}#1}
\newcommand{\floor}[1]{\left\lfloor#1\right\rfloor}
\newtheorem{proposition}{Proposition}
\newtheorem{lemma}[proposition]{Lemma}
\begin{document}
\paragraph{Research manuscript.} OpenAI \texttt{gpt-6-astra}, reasoning \texttt{ultra}; 11 October 2026.
\paragraph{Preserved source body.} The complete catalogue mathematical text and sources follow with layout and source links expanded, including source qualifications and its historical preparation notes. New research begins at the section headed Research attempt.

\section{A scope-locked robust decision certificate for Space Math}\label{statement}

\subsection{Definitions and mathematical statement}

\paragraph{Shared notation.}Standard number systems are $\mathbb N,\mathbb Z,\mathbb Q,\mathbb R,\mathbb C$. Specialized notation and index ranges follow the cited source and the definitions below.


Completing a software project requires an agreed feature set, correctness expectations, documentation, and a release criterion. A decision to continue development can also depend on maintenance resources and intended users. The proper name SpaceMath alone does not define a mathematical object or algorithm.
\paragraph{Record scope and qualifications.}
The catalogue title is retained as metadata. Named certificates, contracts, bounds or frameworks in that title are not asserted as results of the original source extract.
No specification for SpaceMath or criterion for completion is included. The record is a project-management decision, not a mathematical conjecture.

Catalogue screening limitation: Workshop planning or software infrastructure agenda; the supplied passage does not identify a determinate mathematical target.

The source snapshot labels this record partially solved; this is catalogue metadata rather than independent verification of a mathematical result or the current completion of a project.
\paragraph{Statement recorded in the supplied source.}
The sole object of this group is to decide whether or not to finish the software SpaceMath. \href{http://aimpl.org/cyberinfrastructure/7/}{[S2]}
\paragraph{Source status.}
The supplied catalogue labels this source record partially solved; that metadata was checked against the supplied snapshot on 2026-10-11. The label is preserved for provenance and is not an independent verification of present open status or mathematical correctness.
\subsection{Short English statement}
Decide whether to complete the software project called SpaceMath.
\subsection{Sources}
\begin{itemize}
\item[S1] Ulam AI and the UnsolvedMath Contributors, UnsolvedMath: A Curated Collection of Open Mathematics Problems, catalogue entry. Catalogue curation and metadata: CC BY 4.0. The source extract is retained; definitions, English summaries and qualifications are editorial additions.
\url{https://huggingface.co/datasets/ulamai/UnsolvedMath}.
\item[S2] AIM Workshop Problem Lists, source linked by the catalogue.
\url{http://aimpl.org/cyberinfrastructure/7/}.
\end{itemize}
\paragraph{Preparation and review.}
The original source and UnsolvedMath contributors retain attribution for the source extract. Codex (GPT-6; exact serving revision not exposed) assisted with definitions, English summaries, scope notes and formatting. No new human mathematical authorship or independently verified solution is asserted. Catalogue identity checked against the complete supplied registry and current repository. Distinct upstream records, including repeated questions and broad or damaged extracts, are retained for complete coverage of the source CS category. A repeated record is not a new independent conjecture. Complete mathematical-equivalence and current-literature review has not been performed. Preparation and final review follow contribution prompts C and D, with source gaps disclosed. Structural and rendering checks do not establish a new theorem or open status.
\label{end:statement}

\section{Research attempt}
\subsection{Literature, source context and scope}
The complete catalogue body above retains the original agenda, source status, and historical preparation. The AIMPL problem page could not be retrieved. The supplied complete source transcription was read, and the accessible primary AIM workshop announcement was checked for context. The announcement is not a substitute for verification of the inaccessible problem page. The catalogue's earlier proposed frameworks and research summaries are context, not established solutions credited to this attempt. The following is an explicitly limited investigation of one route; it neither substitutes a formalized toy model for the entire source agenda nor claims empirical or deployed-software results. Elementary auxiliary facts are not new original-target progress.

\subsection{A bounded decision calculation, not a project verdict}
The original request is whether to finish Space Math. Its primary introduction was read: it motivates semantic distinctions in mathematical notation and describes a grammar/parser/output route. It does not provide the remaining development cost, success probability, or a decision-maker's valuation. We therefore test a limited numerical decision model without asserting that its numbers describe this software.
Suppose finishing costs $c=4$ utility units and yields benefit $b$, equally likely to be 0 or 10; stopping has utility zero. There are no other actions or costs. A perfect study reveals $b$ before the finish/stop decision.
\begin{proposition}
In this model the optimal expected utility without information is 1, with perfect information is 3, and the maximum expected utility gain from any cost-free information is 2.
\end{proposition}
\begin{proof}
Finishing without information gives expected utility $(0-4+10-4)/2=1$, while stopping gives zero. A randomized mix is a convex combination and cannot improve on 1. With perfect information, stop at $b=0$ and finish at $b=10$, giving $(0+6)/2=3$. Under any possibly randomized signal and policy, the realized action's utility is at most $\max(0,b-4)$ in each state; taking expectation bounds its expected utility by 3. Hence its improvement over the no-information optimum is at most $3-1=2$, attained by the perfect signal.
\end{proof}
The attempted route of deciding from the stated project objective alone stops at missing utilities, costs and uncertainty. In the same two-state model with a different prior probability $p$ of benefit 10, finishing has expected value $10p-4$ and becomes preferable exactly when $p>2/5$; the public description does not determine $p$. Even this threshold presupposes that the two options and utility scale capture the relevant preferences.
\subsection{Verification and remaining gaps}
The two-state expected utilities and the threshold were computed exactly. No project estimate, stakeholder preference, completion decision or current implementation audit was obtained. The calculation is standard value-of-information analysis under invented, clearly stated inputs; it supplies no actual verdict about Space Math.

\section{References and source attribution}
Primary source: \url{http://aimpl.org/cyberinfrastructure/7/}. Attempted primary AIMPL page retrieval; failed with timeout or gateway error. The supplied full source transcription was read; the inaccessible live page was not independently inspected.

\url{https://aimath.org/pastworkshops/cyberinfrastructure.html}. Read the primary December 2023 workshop announcement; it describes the open-source mathematical infrastructure agenda.

\url{https://aimath.org/~farmer/spacemath/}. Read the primary Space Math introduction and examples: semantic ambiguity motivates a formal grammar and conversion to display formats. No budget or completion decision was inferred.

All retrievals were checked on 11 October 2026. Original contributor and catalogue attribution remains in the preserved Sources section. The elementary arguments below are self-contained; their exposition is not a claim of first discovery.


\paragraph{Integration review.} Codex (GPT-6; exact serving revision not exposed) reviewed the written arguments and expanded layout and source-reference wrappers. Repository metadata and the production web renderer were checked using Python 3.13 and JavaScript with MathJax 3.2.2. This model-assisted review is not independent human proof verification; no PDF compilation or formal proof checking is asserted.

\section{Completion Estimate}
\noindent\textbf{COMPLETION: 0\%}
No substantiated new progress toward the entire original target is claimed. The worked results and failed route are delimited above; a known result, illustrative model, or auxiliary obstruction is not counted as original-target completion.

\section{Final status and attribution}
\textbf{Attempt status: unresolved.} This describes this research attempt, not a new ruling on the source catalogue's status or on external literature claims.
The mathematics and this self-audit were generated by OpenAI \texttt{gpt-6-astra}, reasoning \texttt{ultra}, on 11 October 2026. Codex (GPT-6; exact serving revision not exposed) orchestrated the batch and contributed formatting/repository review. No human mathematical authorship, independent human proof checking, or formal verification is claimed. The historical preparation paragraphs in the reproduced catalogue body concern that body only.

Contribution prompt A was applied to the complete definition and supplied source record, followed by the adversarial checks of prompt B. Local reproduction files are \texttt{batch013/checks.py}, \texttt{computations.json}, and the per-ID audit JSON. They were not public links at preparation time. Numerical checks are finite and do not prove unrestricted results. No LaTeX engine was found locally; compilation was not performed by this research agent.

\end{document}
